\documentclass[10pt,a4paper]{amsart}
\usepackage{amsmath,amssymb,amsthm,graphicx}
% For commutative diagrams we use Xy-pic
\usepackage[all]{xy}

% We use 2cell for 2-commutative diagrams.
\xyoption{2cell}
% Theorem environments.
%
\theoremstyle{plain}
\newtheorem{theorem}[subsection]{Theorem}
\newtheorem{proposition}[subsection]{Proposition}
\newtheorem{lemma}[subsection]{Lemma}

\theoremstyle{definition}
\newtheorem{definition}[subsection]{Definition}
\newtheorem{example}[subsection]{Example}
\newtheorem{exercise}[subsection]{Exercise}
\newtheorem{situation}[subsection]{Situation}

\theoremstyle{remark}
\newtheorem{remark}[subsection]{Remark}
\newtheorem{remarks}[subsection]{Remarks}

\numberwithin{equation}{subsection}

% Macros
%
\def\lim{\mathop{\mathrm{lim}}\nolimits}
\def\colim{\mathop{\mathrm{colim}}\nolimits}
\def\Spec{\mathop{\mathrm{Spec}}}
\def\Hom{\mathop{\mathrm{Hom}}\nolimits}
\def\Ext{\mathop{\mathrm{Ext}}\nolimits}
\def\SheafHom{\mathop{\mathcal{H}\!\mathit{om}}\nolimits}
\def\SheafExt{\mathop{\mathcal{E}\!\mathit{xt}}\nolimits}
\def\Sch{\mathit{Sch}}
\def\Mor{\mathop{\mathrm{Mor}}\nolimits}
\def\Ob{\mathop{\mathrm{Ob}}\nolimits}
\def\Sh{\mathop{\mathit{Sh}}\nolimits}
\def\NL{\mathop{N\!L}\nolimits}
\def\CH{\mathop{\mathrm{CH}}\nolimits}
\def\proetale{{pro\text{-}\acute{e}tale}}
\def\etale{{\acute{e}tale}}
\def\QCoh{\mathit{QCoh}}
\def\Ker{\mathop{\mathrm{Ker}}}
\def\Im{\mathop{\mathrm{Im}}}
\def\Coker{\mathop{\mathrm{Coker}}}
\def\Coim{\mathop{\mathrm{Coim}}}

% Boxtimes
%
\DeclareMathSymbol{\boxtimes}{\mathbin}{AMSa}{"02}

%
% Macros for moduli stacks/spaces
%
\def\QCohstack{\mathcal{QC}\!\mathit{oh}}
\def\Cohstack{\mathcal{C}\!\mathit{oh}}
\def\Spacesstack{\mathcal{S}\!\mathit{paces}}
\def\Quotfunctor{\mathrm{Quot}}
\def\Hilbfunctor{\mathrm{Hilb}}
\def\Curvesstack{\mathcal{C}\!\mathit{urves}}
\def\Polarizedstack{\mathcal{P}\!\mathit{olarized}}
\def\Complexesstack{\mathcal{C}\!\mathit{omplexes}}
% \Pic is the operator that assigns to X its picard group, usage \Pic(X)
% \Picardstack_{X/B} denotes the Picard stack of X over B
% \Picardfunctor_{X/B} denotes the Picard functor of X over B
\def\Pic{\mathop{\mathrm{Pic}}\nolimits}
\def\Picardstack{\mathcal{P}\!\mathit{ic}}
\def\Picardfunctor{\mathrm{Pic}}
\def\Deformationcategory{\mathcal{D}\!\mathit{ef}}
\newenvironment{reference}{}{}
\newenvironment{slogan}{}{}
\newenvironment{history}{}{}
\usepackage{cleveref}
\crefname{originallemma}{original lemma}{original lemmas}
\Crefname{originallemma}{Original lemma}{Original lemmas}
\makeatletter
\newlabel{lemma-divisor-special-fiber}{{0C5Z}{source}{}{sourcepointer-1}{}}
\newlabel{lemma-divisor-special-fiber@cref}{{[originallemma][0][]0C5Z}{[1][0][]source}{}{}{}}
\makeatother
\crefname{originaldefinition}{original definition}{original definitions}
\Crefname{originaldefinition}{Original definition}{Original definitions}
\makeatletter
\newlabel{definition-type-minimal}{{0C7A}{source}{}{sourcepointer-2}{}}
\newlabel{definition-type-minimal@cref}{{[originaldefinition][0][]0C7A}{[1][0][]source}{}{}{}}
\makeatother
\crefname{originalproposition}{original proposition}{original propositions}
\Crefname{originalproposition}{Original proposition}{Original propositions}
\makeatletter
\newlabel{curves-proposition-projective-line}{{0C6U}{source}{}{sourcepointer-3}{}}
\newlabel{curves-proposition-projective-line@cref}{{[originalproposition][0][]0C6U}{[1][0][]source}{}{}{}}
\makeatother
\title{A regular proper model is minimal exactly when its numerical type is minimal}
\author{The Stacks Project Authors}
\date{stacks-project:models.tex:local-snapshot-20261012; GNU FDL 1.2 or later}
\begin{document}
\maketitle
\noindent\textit{Editorial note.} The counted unit is the Stacks Project Lemma lemma-numerical-type-minimal-model at source lines 5276-5284 together with its proof at source lines 5286-5309, a straight adjacent pairing in the chapter file. The chapter is the upstream file models.tex, which is the file name recorded in the source record; the pool directory carries the longer name stacks-models and is not used as the chapter name anywhere in this card.\par\medskip
\noindent\textit{Editorial note.} The retained statement opens with the words In Situation followed by a reference to situation-regular-model, so the situation environment that defines the setting is retained as a context span at source lines 4628-4637. It is context, not a counted unit, and the same context span is retained by the separately counted lemma-intersection-pairing; no counted interval is shared.\par\medskip
\noindent\textit{Editorial note.} This chapter's candidate index records the statement body and the proof body without their closing markers, so the retained spans used here are the full environments, exactly as in the earlier cards of this chapter.\par\medskip
\noindent\textit{Editorial note.} The retained proof quotes the author's definition of a minimal numerical type and one proposition of the Algebraic Curves chapter, and the retained situation quotes one lemma of this chapter; all three are declared as non-load-bearing navigation pointers with their stable Stacks tags looked up in the shared tag index. No citation occurs in the retained spans, so no bibliography entry is supplied.\par\medskip
\noindent\textit{Editorial note.} The author's final sentence uses the fact that the Picard group of the projective line over a field is the integers; that is a standard fact about the projective line and is stated inside the proof rather than quoted, so no pointer is needed for it.\par\medskip
\noindent\textit{Editorial note.} The retained preamble is the chapter preamble's Xy-pic setup at lines 14-18, the theorem declarations and mathematical macros at lines 154-217, and the three no-op wrapper environments.\par\medskip
\noindent\textit{Allowed standard prerequisites (editorial).}\begin{itemize}
\item Minimal models, minimal numerical types and exceptional curves of the first kind. In the regular-model situation the numerical type attached to the model records the multiplicities, the genus of each component, the intersection numbers of pairs of components and the residue field degrees; the type is called minimal when no index has both genus zero and self-intersection equal to the negative of the residue field degree, and a curve is called exceptional of the first kind when it is a smooth rational curve with self-intersection minus one over its field of definition. The author's own definitions are the allowed inputs, stated in the same section and in the Algebraic Curves chapter.. The counted lemma is the bridge between the geometric notion of a minimal model and the combinatorial notion of a minimal numerical type, and it is used whenever the chapter switches between the two.
\end{itemize}\medskip
\ifdefined\hypertarget\hypertarget{sourcepointer-1}{}\fi
\noindent\textit{Source reference.} Original lemma 0C5Z: focused-seeds/AGP/sources/stacks-snapshot-20261011/tags.txt, source lines 15279--15279. Navigation pointer inside the retained situation, which names it when it writes the special fibre as a sum of its components with multiplicities; the counted equivalence does not depend on the multiplicities. The number is the stable Stacks tag of the target, looked up in the shared tag index.\par\medskip
\ifdefined\hypertarget\hypertarget{sourcepointer-2}{}\fi
\noindent\textit{Source reference.} Original definition 0C7A: focused-seeds/AGP/sources/stacks-snapshot-20261011/tags.txt, source lines 15325--15325. Pointer to the author's own definition of a minimal numerical type, quoted inside the retained proof; it is an allowed input of the section and is not reproduced inside the span. The number is the stable Stacks tag of the target, looked up in the shared tag index.\par\medskip
\ifdefined\hypertarget\hypertarget{sourcepointer-3}{}\fi
\noindent\textit{Source reference.} Original proposition 0C6U: focused-seeds/AGP/sources/stacks-snapshot-20261011/tags.txt, source lines 15309--15309. Allowed input of the Algebraic Curves chapter quoted by the retained proof to identify a geometrically rational curve of degree one with the projective line; the pointer is registered so the reference resolves and the cited classification is not reproduced inside the span. The number is the stable Stacks tag of the target, looked up in the shared tag index.\par\medskip
\begin{situation}
\label{situation-regular-model}
Let $R$ be a discrete valuation ring with fraction field $K$,
residue field $k$, and uniformizer $\pi$.
Let $C$ be a smooth projective curve over $K$ with $H^0(C, \mathcal{O}_C) = K$.
Let $X$ be a regular proper model of $C$.
Let $C_1, \ldots, C_n$ be the irreducible components of the special
fibre $X_k$. Write $X_k = \sum m_i C_i$ as in
Lemma \ref{lemma-divisor-special-fiber}.
\end{situation}

\begin{lemma}
\label{lemma-numerical-type-minimal-model}
In Situation \ref{situation-regular-model}. The following
are equivalent
\begin{enumerate}
\item $X$ is a minimal model, and
\item the numerical type associated to $X$ is minimal.
\end{enumerate}
\end{lemma}

\begin{proof}
If the numerical type is minimal, then there is no $i$ with
$g_i = 0$ and $(C_i \cdot C_i) = -[\kappa_i: k]$, see
Definition \ref{definition-type-minimal}.
Certainly, this implies that none of the curves $C_i$
are exceptional curves of the first kind.

\medskip\noindent
Conversely, suppose that the numerical type is not minimal.
Then there exists an $i$ such that $g_i = 0$ and
$(C_i \cdot C_i) = -[\kappa_i: k]$.
We claim this implies that $C_i$ is an exceptional curve
of the first kind. Namely, the invertible sheaf
$\mathcal{O}_X(-C_i)|_{C_i}$ has degree $-(C_i \cdot C_i) = [\kappa_i : k]$
when $C_i$ is viewed as a proper curve over $k$, hence
has degree $1$ when $C_i$ is viewed as a proper curve over $\kappa_i$.
Applying
Algebraic Curves, Proposition \ref{curves-proposition-projective-line}
we conclude that $C_i \cong \mathbf{P}^1_{\kappa_i}$ as schemes
over $\kappa_i$. Since the Picard group of $\mathbf{P}^1$
over a field is $\mathbf{Z}$, we see that the normal sheaf
of $C_i$ in $X$ is isomorphic to $\mathcal{O}_{\mathbf{P}_{\kappa_i}}(-1)$
and the proof is complete.
\end{proof}

\end{document}
